% Part: history
% Chapter: biographies
% Section: gerhard-gentzen
%READERNOTE{SOL6-A273: गेंटझेन यांच्या तार्किक निगमनावरील I–II लेखांसाठी Springerची मूळ प्रकाशकनोंद 21 जुलै 1933 हा प्राप्तिदिनांक आणि 1935 हे प्रकाशनवर्ष देते. म्हणून 1934 शी जोडलेले निर्मितीवर्ष सांगण्याऐवजी निश्चित प्रकाशनवर्ष दिले आहे. मूळ लेख: https://doi.org/10.1007/BF01201353 आणि https://doi.org/10.1007/BF01201363. येथे संपूर्ण चरित्राचे स्वतंत्र ऐतिहासिक प्रमाणीकरण केलेले नाही.}

\documentclass[../../../include/open-logic-section]{subfiles}

\begin{document}

\olfileid{his}{bio}{gen}

\olsection{गरहार्ड गेंटझेन}

\olphoto{gentzen-gerhard}{गरहार्ड गेंटझेन}

गरहार्ड गेंटझेन हे मुख्यतः संरचनात्मक सिद्धता उपपत्तीचे निर्माते म्हणून, आणि विशेषतः नैसर्गिक निगमन व क्रमवर्ती कलन या !!{derivation} पद्धती निर्माण केल्याबद्दल ओळखले जातात. त्यांचा जन्म २४ नोव्हेंबर १९०९ रोजी जर्मनीतील ग्राइफ्सवाल्ड येथे झाला. महाविद्यालयपूर्व शाळेत जाण्यापूर्वी गरहार्ड यांचे तीन वर्षे घरीच शिक्षण झाले. शाळेत गेल्यावर शिक्षणाच्या बाबतीत ते आपल्या बहुतेक सहाध्यायांच्या मागे होते. तरीही ते उत्तम विद्यार्थी होते आणि गणितातील त्यांची विशेष क्षमता दिसून येत होती. त्यांची आवड विविध विषयांत होती. उदाहरणार्थ, त्यांनी आपल्या आईसाठी कविता आणि शाळेच्या नाट्यगृहासाठी नाटकेही लिहिली.

गेंटझेन यांनी विद्यापीठीय शिक्षणाची सुरुवात ग्राइफ्सवाल्ड विद्यापीठात केली; त्यानंतर ते गॉटिंगेन (G\"{o}ttingen), म्यूनिख आणि बर्लिन येथे गेले. १९३३ मध्ये त्यांनी गॉटिंगेन (G\"{o}ttingen) विद्यापीठातून हेरमान वाइल यांच्या मार्गदर्शनाखाली डॉक्टरेट मिळवली. (त्यांच्या बहुतांश कामाचे मार्गदर्शन पॉल बर्नायस यांनी केले होते; मात्र नाझींनी त्यांना विद्यापीठातून काढून टाकले.) १९३४ मध्ये गेंटझेन डेव्हिड हिल्बर्ट यांचे सहाय्यक म्हणून काम करू लागले. क्रमवर्ती कलन आणि नैसर्गिक निगमन या !!{derivation} पद्धतींवरील त्यांचे \emph{Untersuchungen \"{u}ber das logische Schlie\ss en I--II
  [Investigations Into Logical Deduction I--II]} हे लेख १९३५ मध्ये प्रकाशित झाले. १९३६ मध्ये त्यांनी पेआनो स्वयंसिद्धकांची सुसंगती सिद्ध केली.

गेंटझेन यांचा नाझींशी असलेला संबंध गुंतागुंतीचा होता. त्यांचे मार्गदर्शक बर्नायस यांना जर्मनी सोडण्यास भाग पाडले जात असतानाच गेंटझेन नाझींची निमलष्करी संघटना असलेल्या एसएच्या विद्यापीठीय शाखेत सामील झाले. अनेक जर्मनांप्रमाणे ते नाझी पक्षाचेही सदस्य होते. युद्धकाळात त्यांनी हवाई गुप्तचर विभागात दूरसंचार अधिकारी म्हणून काम केले. मात्र १९४२ मध्ये मानसिक आरोग्य अचानक खालावल्याने त्यांना सेवेतून मुक्त करण्यात आले. गेंटझेन यांची निष्ठा नाझी पक्षाशी होती की शैक्षणिक प्रगती सुनिश्चित करण्यासाठी ते पक्षात सामील झाले, हे स्पष्ट नाही.

१९४३ मध्ये गेंटझेन यांना प्राग येथील जर्मन विद्यापीठाच्या गणित संस्थेत शैक्षणिक पद देऊ करण्यात आले; त्यांनी ते स्वीकारले. मात्र १९४५ मध्ये प्रागच्या नागरिकांनी जर्मन अधिपत्याविरुद्ध बंड केले. सोव्हिएत सैन्य शहरात दाखल झाले आणि विद्यापीठातील सर्व प्राध्यापकांना अटक केली, असे या स्रोतामध्ये म्हटले आहे. नाझी संघटनांतील सदस्यत्वामुळे गेंटझेन यांना सक्तीच्या श्रमांच्या छावणीत नेण्यात आले. ४ ऑगस्ट १९४५ रोजी वयाच्या ३५व्या वर्षी कारागृहातील त्यांच्या कोठडीत कुपोषणामुळे त्यांचे निधन झाले.

\begin{reading}
गेंटझेन यांचे संपूर्ण चरित्र \citet{Menzler-Trott2007} येथे पाहा. नाझी राजवटीखालील गणितज्ञांविषयीचे \citet{Segal2014} यांचे पुस्तकही वाचनीय आहे; त्यात गेंटझेन यांच्या आयुष्याविषयी एक संक्षिप्त नोंद आहे. तार्किक निगमनावरील गेंटझेन यांचे मूळ जर्मन लेख \citep{Gentzen1935a,Gentzen1935b} येथे उपलब्ध आहेत. त्यांच्या लेखांची इंग्रजी भाषांतरे \citet{Gentzen1969} यांनी एका खंडात एकत्र केली आहेत; त्यात त्यांचे संक्षिप्त चरित्रही आहे.
\end{reading}

\end{document}
